% ECONOMICS_PROBLEM: {"id": "D44-127", "title": "Constant prophet bounds for subadditive bidders: corrected target", "jel_code": "D44", "source_id": "OP-0175"}
% !TeX program = xelatex
\documentclass[11pt,a4paper]{article}
\usepackage[margin=23mm]{geometry}
\usepackage{fontspec}
\setmainfont{Latin Modern Roman}
\usepackage{amsmath,amssymb,mathtools}
\usepackage{xcolor,enumitem,booktabs,longtable,array,xurl,hyperref}
\definecolor{muted}{HTML}{53636C}
\hypersetup{colorlinks=true,urlcolor=blue,linkcolor=blue}
\setlength{\parindent}{0pt}
\setlength{\parskip}{5pt}
\newcommand{\R}{\mathbb R}
\newcommand{\E}{\mathbb E}
\newcommand{\N}{\mathbb N}
\newcommand{\ind}{\mathbf 1}
\newcommand{\Var}{\operatorname{Var}}
\newcommand{\Cov}{\operatorname{Cov}}
\newcommand{\supp}{\operatorname{supp}}
\DeclareMathOperator*{\argmin}{arg\,min}
\DeclareMathOperator*{\argmax}{arg\,max}
\newcommand{\entrymeta}[3]{{\sffamily\small\color{muted}\textbf{\texttt{#1}}\quad|\quad #2\par #3\par}}
\newcommand{\Problemref}[1]{\texttt{#1}}
\newcommand{\JELref}[2]{\texttt{#2} (JEL \texttt{#1})}
\begin{document}
\section*{Constant prophet bounds for subadditive bidders: corrected target}
\textbf{Problem ID:} \texttt{D44-127}\quad \textbf{JEL:} \texttt{D44}\par
% BEGIN ECONOMICS STATEMENT
\label{problem:OP-0175}
% GAME_THEORY_PROBLEM: {"local_id": "GT-045", "original_number": 45, "kind": "P", "entry_kind": "statement_only", "published": false, "review_required": true, "category": "game_theory"}
\label{game:GT-045}
{\sffamily\small\color{muted}
\textbf{GT-045} | Online mechanisms | \textbf{P}: posted-item-price question; unrestricted prophet question resolved.
\par}
\subsection*{Definitions and mathematical statement}
There are $m$ items and $n$ bidders arriving in order $1,\ldots,n$. Their independent random normalized monotone valuations satisfy $v_i(S\cup T)\leq v_i(S)+v_i(T)$. The distributions are known, and $\mathbb E\operatorname{OPT}(v)<\infty$. A static posted-item-price mechanism draws $p\in\mathbb R_{\geq0}^m$ independently of realized values before arrivals. At arrival $i$, bidder $i$ purchases a utility-maximizing subset $S_i$ of remaining items at payment $\sum_{g\in S_i}p_g$. Tie-breaking is specified as part of the mechanism and selects only maximizing bundles. The corrected question is whether an absolute constant $c>0$ exists such that for every instance some distribution over prices satisfies
\[
 \mathbb E_{v,p}\Big[\sum_i v_i(S_i)\Big]\geq c\,\mathbb E_v[\operatorname{OPT}(v)].
\]
This is an existence/competitive-ratio question; efficient price computation is a separate requirement.
\subsection*{Short English statement}
Can static item prices obtain a constant fraction of offline welfare for independent subadditive bidders?
\subsection*{Variants and scope boundaries}
The unconstrained online allocation problem already admits a constant bound. Its algorithm need not be implementable by item prices or DSIC. Bundle prices and adaptive prices enlarge the present class.
\subsection*{Source relationship and status}
The original catalogue incorrectly retained the general prophet-existence question: Correa and Cristi establish a $1/6$ bound in 2023. Their posted-prices discussion keeps the more restrictive mechanism question separate and suggests a positive answer. This corrected statement follows that discussion, not an assertion that the general bound is unknown.
\subsection*{References}
\begingroup\footnotesize\raggedright
\textbf{[S1]} Jos\'e Correa and Andr\'es Cristi (2023). \emph{A Constant Factor Prophet Inequality for Online Combinatorial Auctions}. STOC 2023, pp. 686--697. \textit{Verified locator:} introduction, $1/6$ result; manuscript printed pp. 14--16, ``Posted Prices.'' \url{https://www.dii.uchile.cl/~jcorrea/papers/Manuscripts/2023CC.pdf}
\par\endgroup

\subsection*{Common conventions: Source mathematical conventions}

\textbf{Finite games.} For a finite set $A$, $\Delta(A)$ is its probability
simplex. Players choose independent mixed strategies in Nash equilibrium;
correlation is allowed only when explicitly part of the solution concept.
In a game with payoff functions $u_i$ and strategy profile $x$, an additive
$\varepsilon$ Nash equilibrium satisfies
\[
 u_i(x)\geq u_i(x_i',x_{-i})-\varepsilon
 \quad\text{for every player $i$ and every admissible deviation $x_i'$}.
\]
For numerical approximation, payoffs are normalized as locally specified.
An $\varepsilon$ payoff residual is distinct from distance at most
$\varepsilon$ to the set of exact equilibria. Point convergence, convergence
of empirical play, and convergence in distance to an equilibrium set are
also distinct.

\textbf{Computational representations.} Unless stated otherwise, a complexity
question takes rational data with binary encoding; polynomial time is measured
in total bit length. A fixed-accuracy polynomial algorithm is not necessarily
polynomial in $1/\varepsilon$, and neither is necessarily polynomial in
$\log(1/\varepsilon)$. Succinct games and value, demand, or sampling oracles
must declare their interfaces. Exact real solutions require an explicit
representation; an unspecified list of arbitrary real numbers is not a
finite computational output. A claimed algorithmic barrier is meaningful
only for its specified model and reductions.

\textbf{Dynamic games.} Histories, information, observation, and allowable
randomization are part of the model. A stationary strategy depends only on
the current state; a positional strategy in a deterministic graph game
chooses one action at each controlled vertex. Nash equilibrium at an initial
state and equilibrium after every history are different requirements.
An expectation of a pathwise $\liminf$ payoff, a $\liminf$ of expected averages,
a limiting expected average, and a uniform finite-horizon equilibrium are
different criteria. None is silently substituted for another.

\textbf{Allocation and mechanisms.} A complete allocation assigns every item
exactly once unless otherwise stated. Zero-valued goods, negative values,
capacity constraints, feasibility, lotteries, and copies of goods affect
fairness definitions and are specified locally. Dominant-strategy incentive
compatibility, Bayesian incentive compatibility, universal truthfulness,
and truthfulness in expectation impose different inequalities. Whenever
an objective ratio has zero denominator, its convention is stated or those
instances are excluded.

\textbf{Infinite-dimensional models.} Expectations, integrals, stochastic
equations, and optimization objectives require their local regularity,
integrability, filtrations, and admissible domains. A backward--forward
mean-field system is a boundary-value problem; it is not an initial-value
semiflow without an additional construction. Long-time, large-population,
and vanishing-noise limits require an order of limits and a convergence
topology. If a minimum need not be attained, the objective is its infimum
or supremum and the approximation requirement is stated separately.
% END ECONOMICS STATEMENT
\end{document}
